\documentclass[10pt,letterpaper]{article}
\usepackage{cornell-notes}

\title{Clause 26.07.08.04: Class Filter View Sentinel}
\author{}
\date{}
\setDocTitle{Clause 26.07.08.04: Class Filter View Sentinel}
\setDocSubtitle{C++ 2024}
\setDocOwner{C++ 2024 Cornell Notes}
\setDocDate{\today}
\setCornellCollection{C++ 2024 Cornell Notes}
\setCornellUnitType{Clause}
\setCornellUnitNumber{26.07.08.04}
\setCornellUnitTitle{Class Filter View Sentinel}
\hypersetup{pdftitle={Clause 26.07.08.04: Class Filter View Sentinel},pdfsubject={Cornell notes},pdfauthor={}}

\begin{document}
\maketitle
\begin{CornellOverviewBox}{Overview}
\textbf{Classification:} Standard library - Normative\\[2pt]
\textbf{Learning objective:} Understand the rules, terminology, and semantic consequences associated with Class filter\_view::sentinel.
\end{CornellOverviewBox}\begin{CornellSummaryBox}{Summary}
{\sffamily\bfseries\color{Accent} Summary}\par\vspace{3pt}Section 26.7.8.4 focuses on Class filter\_view::sentinel. Apply the specified interface together with its mandates, effects, result conditions, exception behavior, and complexity constraints. When applying it, preserve the distinction between mandatory requirements, recommendations, permissions, and behavior deliberately left undefined, unspecified, or implementation-defined.
\end{CornellSummaryBox}

\end{document}
